\documentclass[12pt]{article}
\usepackage[T1]{fontenc}
\usepackage[utf8]{inputenc}
\usepackage{lmodern}
\usepackage{textcomp}
\usepackage{enumitem}
\usepackage{geometry}
\geometry{margin=1in}
\begin{document}
\begin{center}
Answer Key - Chemistry - Undergraduate Level Assessment 24 \
\small (Answers are ordered exactly as in the question paper.)
\end{center}

\section*{Multiple Choice}
\begin{enumerate}[label=\arabic*.]
\item \textbf{Answer:} D
\\ \textit{Options:} A) Melting point; B) Boiling point; C) Polarizability; D) First ionization energy
\\ \textit{Note:} The first ionization energy of argon is lower than that of neon because argon has a larger atomic radius and more electron shielding, which reduces the effective nuclear charge experienced by its outermost electrons, making them easier to remove.
\item \textbf{Answer:} B
\\ \textit{Options:} A) The I- : IO 3- ratio is 3:1.; B) The MnO 4- : I- ratio is 6:5.; C) The MnO 4- : Mn$_{2}$+ ratio is 3:1.; D) The H+ : I- ratio is 2:1.
\\ \textit{Note:} The correct answer is based on the stoichiometry of the balanced redox reaction, which shows that for every 6 moles of permanganate ions (MnO 4-) reacting, 5 moles of iodide ions (I-) are consumed, resulting in a 6:5 ratio.
\item \textbf{Answer:} B
\\ \textit{Options:} A) 23.56 GHz; B) 42.58 MHz; C) 74.34 kHz; D) 13.93 MHz
\\ \textit{Note:} The Larmor frequency for a proton in a magnetic field of 1 T is calculated using the formula \( f = \gamma B \), where \( \gamma \) (the gyromagnetic ratio for a proton) is approximately 42.58 MHz/T, thus yielding a frequency of 42.58 MHz when the magnetic field strength is 1 T.
\item \textbf{Answer:} A
\\ \textit{Options:} A) HCl; B) AgCl; C) CaCl 2; D) CCl 4
\\ \textit{Note:} HCl has the lowest melting point among the given compounds due to its weak intermolecular forces, specifically dipole-dipole interactions, which are less substantial than the stronger forces present in other compounds.
\item \textbf{Answer:} D
\\ \textit{Options:} A) BeCl 2; B) MgCl 2; C) ZnCl 2; D) SCl 2
\\ \textit{Note:} SCl 2 is least likely to behave as a Lewis acid because it has a complete octet and a stable electronic configuration, making it less inclined to accept an electron pair compared to other compounds that are more electron-deficient.
\end{enumerate}

\section*{True / False}
\begin{enumerate}[label=\arabic*.]
\setcounter{enumi}{5}
\item \textbf{Answer:} True
\\ \textit{Options:} A) True; B) False
\\ \textit{Note:} Most organic compounds are primarily composed of carbon, hydrogen, oxygen, and nitrogen, and since silicon is not a common element in organic chemistry, tetramethylsilane, which contains silicon, provides a unique and non-polar reference for 1 H NMR chemical shifts.
\item \textbf{Answer:} False
\\ \textit{Options:} A) True; B) False
\\ \textit{Note:} The statement is false because the hybrid sp³ orbitals for carbon are constructed from a linear combination of only three atomic orbitals: one 2 s orbital and three 2 p orbitals, not four.
\item \textbf{Answer:} False
\\ \textit{Options:} A) True; B) False
\\ \textit{Note:} The statement is false because although hexane is symmetrical, it has multiple distinct carbon environments, leading to more than five distinct chemical shifts in its 13 C NMR spectrum.
\item \textbf{Answer:} True
\\ \textit{Options:} A) True; B) False
\\ \textit{Note:} Protons in a molecule with a rotational correlation time of 1 ns exhibit rapid spin-lattice relaxation at a magnetic field of 3.74 T due to their optimal interaction with the magnetic field, allowing for efficient energy exchange between the spins and the lattice.
\item \textbf{Answer:} False
\\ \textit{Options:} A) True; B) False
\\ \textit{Note:} The statement is false because ionic strength is determined not only by the concentration of ions but also by their charge; 0.050 M AlCl 3 produces more ions (3 Cl⁻ and 1 Al³⁺) than 0.100 M NaCl (2 Na⁺ and 2 Cl⁻), resulting in a higher ionic strength for AlCl 3 despite its lower molarity.
\end{enumerate}

\section*{Long Form Response}
\begin{enumerate}[label=\arabic*.]
\setcounter{enumi}{10}
\item \textbf{Answer:} Electron affinity refers to the energy change that occurs when an electron is added to a neutral atom, and several factors influence this property, including atomic size, nuclear charge, and electron configuration. Generally, elements with a higher effective nuclear charge and smaller atomic radius exhibit greater electron affinity due to the stronger attractive forces on the added electron. Calcium, located in Group 2 of the periodic table, has the lowest electron affinity among its peers primarily because it has a relatively stable electron configuration with a full outer s subshell, making it less favorable to gain an additional electron. Furthermore, the larger atomic size of calcium compared to nonmetals results in a weaker attraction for the added electron, contributing to its lower electron affinity.
\\ \textit{Note:} The answer correctly identifies that electron affinity is influenced by factors such as atomic size and nuclear charge, and it specifically notes that calcium's position in Group 2, with its filled s-orbital, results in a lower tendency to gain an electron compared to other elements, leading to the conclusion that it has the lowest electron affinity among its peers.
\item \textbf{Answer:} In chemistry, conjugate acid-base pairs consist of two species that can transform into each other by the gain or loss of a proton (H+). The pair of H$_{3}$O+ (hydronium ion) and H$_{2}$O (water) exemplifies this concept, where H$_{3}$O+ serves as the conjugate acid of H$_{2}$O, which acts as the conjugate base. In an acid-base reaction, H$_{2}$O can accept a proton to become H$_{3}$O+, indicating its role as a base, while H$_{3}$O+ can donate a proton back to revert to H$_{2}$O, demonstrating its function as an acid. These interactions are fundamental in acid-base chemistry, illustrating how the transfer of protons facilitates various chemical processes in aqueous solutions.
\\ \textit{Note:} The answer correctly identifies H$_{3}$O+ and H$_{2}$O as a conjugate acid-base pair, illustrating their roles in acid-base reactions through the processes of proton donation and acceptance, which exemplify the fundamental concept of conjugate acids and bases.
\end{enumerate}

\vfill
\noindent\textit{End of Answer Key}
\end{document}